% =====================================================================
%  ch07_uart_rx.tex
% =====================================================================
\section{串口接收模块 \texttt{uart\_rx}（\file{src/uart_rx.v}）}

\subsection{文件信息}

\begin{center}
\small
\begin{tabular}{ll}
\toprule
文件路径 & \file{FPGA/src/uart_rx.v}（注释后 187 行）\\
模块 / 实例 & \texttt{uart\_rx} / \sig{uart_rx_inst}\\
帧格式 & 8 数据位、1 停止位、无校验（8N1），LSB 先到\\
参数 & \sig{UART_BPS}=115200，\sig{CLK_FREQ}=50\,000\,000\\
派生参数 & \sig{BAUD_CNT_MAX}=434\\
\bottomrule
\end{tabular}
\end{center}

\subsection{第 1 层：接口与参数}

\begin{table}[htbp]
\centering
\small
\begin{tabular}{llll}
\toprule
端口 & 方向 & 位宽 & 说明\\
\midrule
\sig{sys_clk} & input & 1 & 50MHz\\
\sig{sys_rst_n} & input & 1 & 异步复位\\
\sig{rx} & input & 1 & 串口输入（异步信号）\\
\sig{po_data} & output & 8 & 收到的字节\\
\sig{po_flag} & output & 1 & 字节有效脉冲，宽 1 个时钟\\
\bottomrule
\end{tabular}
\end{table}

分频参数计算：
\begin{equation}
\text{BAUD\_CNT\_MAX}=\left\lfloor\frac{50\,000\,000}{115\,200}\right\rfloor=434,
\qquad
\text{实际波特率}=\frac{50\,000\,000}{434}=115\,207\ \text{bps},
\end{equation}
误差约 $+0.006\%$，远小于 UART 允许的 $\pm2\%$。

\subsection{第 3 层：语句块解析}

\subsubsection{块 1--3：三级同步器（57--82 行）}

\begin{lstlisting}[style=vsnip]
always@(posedge sys_clk or negedge sys_rst_n)
    if(sys_rst_n==1'b0)
        rx_reg1<=1'b1;
    else
        rx_reg1<=rx;
\end{lstlisting}

三个完全相同的寄存器串联：\sig{rx}$\rightarrow$\sig{rx_reg1}$\rightarrow$\sig{rx_reg2}
$\rightarrow$\sig{rx_reg3}。作用是让异步的串口输入在本时钟域内“稳定”下来，
消除亚稳态。串口空闲为高电平，因此复位值取 1。

\subsubsection{块 4：起始位检测（84--94 行）}

\begin{lstlisting}[style=vsnip]
always@(posedge sys_clk or negedge sys_rst_n)
    if(sys_rst_n==1'b0)
        start_flag<=1'b0;
    else if((rx_reg3==1'b1)&&(rx_reg2==1'b0)&&(work_en==1'b0))
        start_flag<=1'b1;
    else
        start_flag<=1'b0;
\end{lstlisting}

在同步后的信号上检测\textbf{下降沿}（前两拍为 1、上一拍为 0），且当前空闲
（\sig{work_en}=0）时产生单周期起始标志。用 \sig{rx_reg2}/\sig{rx_reg3}
错开一拍判断，可避免把抖动误判为起始位。

\subsubsection{块 5：接收忙标志（96--108 行）}

\begin{lstlisting}[style=vsnip]
always@(posedge sys_clk or negedge sys_rst_n)
    if(sys_rst_n==1'b0)
        work_en<=1'b0;
    else if(start_flag==1'b1)
        work_en<=1'b1;
    else if ((bit_cnt==5'd8)&&(bit_flag==1'b1))
        work_en<=1'b0;
    else
        work_en<=work_en;
\end{lstlisting}

检出起始位后置忙，收满 8 个数据位（\sig{bit_cnt}=8）后释放。忙标志有两个作用：
禁止在接收过程中重复触发起始检测；作为 \sig{baud_cnt} 的清零条件。

\subsubsection{块 6/7：位内计数与中点采样（110--131 行）}

\begin{lstlisting}[style=vsnip]
always@(posedge sys_clk or negedge sys_rst_n)
    if(sys_rst_n==1'b0)
        baud_cnt<=16'd0;
    else if((baud_cnt==BAUD_CNT_MAX-1)||(work_en==1'b0))
        baud_cnt<=16'd0;
    else
        baud_cnt<=baud_cnt+16'd1;

always@(posedge sys_clk or negedge sys_rst_n)
    if(sys_rst_n==1'b0)
        bit_flag<=1'b0;
    else if(baud_cnt==BAUD_CNT_MAX/2-1)
        bit_flag<=1'b1;
    else
        bit_flag<=1'b0;
\end{lstlisting}

\sig{baud_cnt} 在 0$\sim$433 循环；当它等于 $434/2-1=216$ 时产生 \sig{bit_flag}，
即在每个位的**正中间**采样——这是 UART 接收抗时钟偏差的标准做法。

\begin{table}[htbp]
\centering
\small
\caption{接收时序}
\begin{tabular}{lll}
\toprule
量 & 公式 & 数值\\
\midrule
位周期 & 434 个时钟 & 8.68 $\mu$s\\
采样时刻 & 第 216 个时钟 & 位中点\\
一帧（起始+8 数据+停止） & 10 位 & 86.8 $\mu$s\\
\bottomrule
\end{tabular}
\end{table}

\subsubsection{块 8/9：位计数与数据移位（133--152 行）}

\begin{lstlisting}[style=vsnip]
always@(posedge sys_clk or negedge sys_rst_n)
    if(sys_rst_n==1'b0)
        bit_cnt<=5'b0;
    else if((bit_cnt==5'd8)&&(bit_flag==1'b1))
        bit_cnt<=5'b0;
    else if(bit_flag==1'b1)
        bit_cnt<=bit_cnt+1'b1;

always@(posedge sys_clk or negedge sys_rst_n)
    if(sys_rst_n==1'b0)
        rx_data<=8'b0;
    else if(( bit_cnt<=5'd8)&&(bit_cnt>=5'd1)&&(bit_flag==1'b1))
        rx_data<={rx_reg3,rx_data[7:1]};
\end{lstlisting}

\sig{bit_cnt}=0 对应起始位中点（不采样），1$\sim$8 分别对应 $d_0\sim d_7$ 的中点。
每次采样把 \sig{rx_reg3} 塞到最高位、整体右移一位，因此先到的 $d_0$ 最终落在 bit0，
符合 UART 的 LSB-first 约定。

\subsubsection{块 10--12：完成脉冲与输出（154--184 行）}

\begin{lstlisting}[style=vsnip]
always@(posedge sys_clk or negedge sys_rst_n)
    if(sys_rst_n==1'b0)
        rx_flag<=1'b0;
    else if((bit_cnt==5'd8)&&(bit_flag==1'b1))
        rx_flag<=1'b1;
    else
        rx_flag<=1'b0;

always@(posedge sys_clk or negedge sys_rst_n)
    if(sys_rst_n==1'b0)
        po_data<=8'b0;
    else if(rx_flag==1'b1)
        po_data<=rx_data;

always@(posedge sys_clk or negedge sys_rst_n)
    if(sys_rst_n==1'b0)
        po_flag<=1'b0;
    else
        po_flag<=rx_flag;
\end{lstlisting}

第 8 个数据位采完后产生 \sig{rx_flag}，把移位结果锁存到 \sig{po_data}，再打一拍得到
\sig{po_flag}，使“数据 + 有效”同拍输出给顶层组帧逻辑。

\begin{warnbox}
\textbf{边界条件}：本实现不检查停止位、不做帧错误检测，也不做接收 FIFO。
若上位机连续高速发送（间隔小于 1 个位周期），可能因 \sig{work_en} 释放晚于下一帧起始沿而丢字节。
对“偶发指令帧”的应用（本工程）足够；若要连续数据流收发，应加 FIFO 与帧错误标志。
\end{warnbox}

\subsection{完整注释源码}

\lstinputlisting[style=vcode]{../FPGA/src/uart_rx.v}
